\clearpage\fancyhead[L]{\small POLARSTERN.}\hypertarget{native-398-part-lower}{}\label{native-398-part-lower}
\begin{center}\hrule\vspace{4pt}\fontsize{10}{12}\selectfont
V. \textsc{Lindenau} u. \textsc{Bohnenberger} Zeitschrift für Astronomie u. verwandte Wiss.\par
Band IV. Seite 119..131. Juli u. August. 1817.\par\vspace{4pt}\hrule\vspace{4pt}\end{center}
\begin{flushright}Göttingen, den 18. Juli 1817.\end{flushright}
\par\indent Wenn ich mich nicht irre, mein teuerster Freund, habe ich Ihnen meine Beobachtungen des Polarsterns mit dem zwölfzolligen \textsc{Reichenbach}schen Kreis noch nicht vollständig mitgeteilt. Ich freue mich, dass ich \SourceLeaf{399} sie jetzt mit einer Reihe von Beobachtungen zur Sonnenwende verbinden kann, die in diesem Jahr durch das schönste Wetter mehr als je begünstigt wurden. Diese Beobachtungen wurden an der neuen Sternwarte gemacht. Deren Bau ist, wie Sie wissen, noch nicht ganz vollendet, und der eigentliche Betrieb wird erst beginnen, wenn die neuen Meridianinstrumente eingetroffen und aufgestellt sind. Nach meiner geodätischen Bestimmung liegt die Mitte der neuen Sternwarte $6''11$ südlich und $1^{\mathrm s}90$ östlich von der Mitte der alten. Der Platz, an dem ich den Polarstern beobachtet habe, liegt $0''12$ nördlich, der Platz der Sonnenbeobachtungen $0''16$ südlich von der Mitte.

\par\indent Sie erinnern sich, dass ich in einigen meiner früheren Briefe eine Vermutung über den viel besprochenen Unterschied zwischen der aus Beobachtungen der Zirkumpolarsterne und der aus Sonnenbeobachtungen bestimmten Polhöhe erwähnt habe. Bei den kleineren Repetitionskreisen könnte dieser Unterschied auf eine Durchbiegung des Fernrohrs auf der Objektivseite zurückgehen, da man am Objektiv einen kräftigen Messingring als Gegengewicht für das Prismenokular und die Klammer aufsteckt. Falls dies eine nachteilige Durchbiegung verursacht, müsste man nach \textsc{Reichenbachs} Rat das aufgesteckte Gewicht so weit verringern, dass es gerade das Fernrohr richtig ausbalanciert. Die Klammer wäre durch ein am Alhidadenkreis angeschraubtes Gewicht auszubalancieren. Für meinen Zweck war es jedoch noch günstiger, das aufgesteckte Gewicht ganz abzunehmen und das vollständige Gleichgewicht durch ein am Alhidadenkreis angeschraubtes Gewicht herzustellen. Wäre jene Vermutung begründet gewesen, so hätte die verbleibende Durchbiegung auf der Okularseite des Fernrohrs nun offensichtlich einen Fehler in der entgegengesetzten Richtung verursachen müssen. Wären zuvor alle Zenitdistanzen infolge einer Durchbiegung zu klein gemessen worden, so müssten sie nun alle zu groß werden, und die Beobachtungen des Polarsterns müssten eine zu kleine, die Sonnenbeobachtungen eine zu große Polhöhe ergeben. Das Ergebnis hat dies jedoch keineswegs bestätigt. Die an der neuen Sternwarte aus dem Polarstern mit angeschraubtem Gewicht bestimmte Polhöhe ist fast haarscharf dieselbe wie diejenige, die sich aus der Übertragung der mit aufgestecktem Gewicht an der alten Sternwarte beobachteten Polhöhe ergibt. Die Sonnenbeobachtungen weichen in derselben Richtung und um denselben Betrag ab wie zuvor. Anfangs, bevor das Wetter zusammenhängende astronomische Beobachtungen zuließ, hatte ich einen anderen Weg gewählt, um die Vermutung zu prüfen. Ich hatte mehrere Reihen von Zenitdistanzen eines Gegenstands auf der Erde beobachtet, einige mit aufgestecktem, die anderen mit angeschraubtem Gewicht. Die ersten Reihen ergaben tatsächlich einen Unterschied in der Richtung und ungefähr in der Größe, die zu erwarten gewesen wären, wenn dies wirklich die Ursache gewesen wäre. In den folgenden Reihen verschwand dieser Unterschied jedoch fast ganz. Überhaupt stimmten die Beobachtungen des irdischen Gegenstands, den man nicht so genau anvisieren konnte wie Gegenstände am Himmel, weniger gut überein, als für diese Prüfung nötig gewesen wäre.

{\rightskip=0pt plus 1em \par\indent Wenn dieser äußerst bemerkenswerte Unterschied auf das Instrument zurückgeht, muss es in seiner Bauweise etwas geben, das dazu führt, dass wir die Zenitdistanzen damit zu klein beobachten. Jenes Gewicht war nicht die Ursache. Nach sorgfältiger Prüfung habe ich auch festgestellt, dass mehrere andere Umstände, die ich in Verdacht zu ziehen geneigt war, keinen Einfluss hatten. Ich gestehe, dass ich bisher keinen Umstand gefunden habe, dem ich eine \emph{solche} Wirkung zuschreiben könnte. Eine Wirkung in \emph{entgegengesetzter} Richtung wäre viel leichter zu erklären und kann tatsächlich eintreten, wenn man nicht auf der Hut ist. Wie Sie wissen, habe ich selbst im vergangenen Winter eine solche Erfahrung gemacht. Ich hatte den Kreis ganz zerlegt, gereinigt und wieder zusammengesetzt. Meine ersten Beobachtungen des Polarsterns ergaben wiederholt eine um $20''$ zu kleine Polhöhe. Es gab also etwas, das die Zenitdistanzen vergrößerte. Nach mehreren anderen Prüfungen fand ich die Fehlerquelle bald in einer Durchbiegung des Okularrohrs, das an seinem äußersten Ende im Fernrohr von der Feder \SourceLeaf{400} festgehalten wird. Die Feder hatte vielleicht etwas an Spannkraft verloren. Nachdem sie durch leichtes Biegen die nötige Spannkraft zurückerhalten hatte, verschwand der Fehler sofort. Derselbe Fehler tritt auch auf, wenn der Faden, an dem die Zenitdistanzen beobachtet werden, nicht richtig gespannt ist. Der Faden hängt dann in der Mitte unter seinem eigenen Gewicht durch, und man beobachtet offensichtlich zu große Zenitdistanzen. Es braucht wohl kaum daran erinnert zu werden, dass es hier um geringe Fehler geht. Wenn der schlaffe Faden aber in der Mitte um ein paar Sekunden durchhängt und seine gesamte Länge vielleicht mehr als ein Grad beträgt, erkennt das Auge noch keine Krümmung, es sei denn, dass man das Bild einer vollkommen geraden waagerechten Linie aus geeigneter Entfernung auf den Faden projizieren lässt.\par}

\par\indent Gewiss ist jedoch, dass man das Instrument nur für fehlerfrei erklären kann, wenn man den entgegengesetzten Fehler anderen Instrumenten zuschreibt, die zwischen den Stern- und Sonnenbeobachtungen entweder keinen solchen Unterschied oder einen kleineren ergeben. Hat der 12-zollige Kreis recht, so geben die dreifüßigen Kreise mit aufrechter Säule und \textsc{Bessels} Kreis zu große Zenitdistanzen. Hat der dreifüßige Kreis recht, so gibt der zwölfzollige zu kleine und der \textsc{Bessel}sche zu große Zenitdistanzen. Hat der letztere recht, so geben alle \textsc{Reichenbach}schen Kreise zu kleine Zenitdistanzen. Mit \emph{Gewissheit} kann man aber keinem Instrument recht geben, solange man nicht entweder die mögliche Fehlerquelle an den anderen Instrumenten eindeutig und unwiderlegbar nachgewiesen oder die Unrichtigkeit der mit den anderen Instrumenten gefundenen Ergebnisse auf anderem Wege \emph{entschieden} dargelegt hat. So wie die Dinge jetzt liegen, müssen wir, scheint mir, noch eingestehen, dass kein bestimmtes Urteil möglich ist. Die relative \emph{Wahrscheinlichkeit} wird jeder Astronom aus subjektiven Gründen vielleicht anders abwägen. Denn ich finde es durchaus natürlich, dass man, wenn sich gegen keines der Instrumente ein bestimmter Einwand erheben lässt, dies am stärksten bei dem Instrument empfindet, mit dem man selbst beobachtet und daher am besten vertraut ist.

\par\indent Zwei weitere Bemerkungen kann ich mir hier jedoch nicht versagen.

\par\indent \emph{Erstens} scheint mir, dass man das Zeugnis der zwölfzolligen Kreise nicht allein \emph{deswegen} verwerfen kann, \emph{weil} sie die kleinsten dieser Instrumente sind, sofern man nicht nachweisen kann, dass sie \emph{wegen ihrer geringen Größe} konstante Fehler verursachen, die immer in dieselbe Richtung wirken. Zufällige Fehler, die bald so, bald anders wirken, könnte man wegen der geringen Größe eher erwarten. Die Erfahrung zeigt aber, dass ihre Ergebnisse ebenso gut oder wenigstens nicht viel schlechter übereinstimmen als die der größeren Instrumente. Ich weiß nicht, ob man \emph{im Allgemeinen} bei größeren Instrumenten wegen ihrer Größe nicht mindestens ebenso sehr vor konstanten Fehlern auf der Hut sein muss.

\par\indent \emph{Zweitens} gestehe ich, dass, wenn man von allen anderen Umständen absieht, bei einem Instrument noch kein Grund für die Annahme seiner Richtigkeit zu bestehen scheint, nur \emph{deswegen}, \emph{weil} es überhaupt keinen Unterschied zwischen Stern- und Sonnenbeobachtungen ergibt. Abgesehen davon, dass noch nicht ganz geklärt ist, ob die Refraktion für die Sonne und für die Fixsterne \emph{genau} dieselbe ist, kommt hier noch ein anderer wichtiger Umstand in Betracht. Wie Sie sich erinnern, habe ich mich darüber schon vor mehreren Jahren mit Ihnen unterhalten und ihn damals als eine mögliche Erklärung des Rätsels angeführt, ohne ihm besonders viel Gewicht beimessen zu wollen. Ich greife den Gegenstand hier wieder auf, da meines Wissens bisher noch nicht öffentlich davon die Rede war. Wir \emph{beobachten} den Mittelpunkt der Gestalt der Sonne. Die Anziehungskraft der Sonne wird dagegen auf die Planeten annähernd so wirken, als wäre sie in ihrem Schwerpunkt vereinigt. Der Schwerpunkt der Sonne wird in der Ebene der Erdbahn liegen, während der Mittelpunkt ihrer Gestalt außerhalb dieser Ebene liegen wird. Man darf nicht annehmen, dass beide mathematisch genau zusammenfallen, da der Sonnenkörper zweifellos keine \emph{ganz} gleichartige Masse ist. Läge der Schwerpunkt etwas südlich vom Mittelpunkt der Gestalt, so hätte dieser immer \SourceLeaf{401} eine kleine nördliche Breite. Soweit man sie außer Acht lässt, müsste sich aus Sonnenbeobachtungen auf der nördlichen Erdhalbkugel immer eine zu kleine Polhöhe ergeben. So einfach und natürlich diese Erklärung auch scheint, sie ist doch mit Schwierigkeiten verbunden. Man müsste erst genauer untersuchen, inwieweit eine solche Zusammensetzung des Sonnenkörpers möglich ist, bei der der Schwerpunkt um $\frac{1}{100}$ oder um $\frac{1}{1000}$ des Sonnendurchmessers südlich vom Mittelpunkt der Gestalt läge, ohne dass die Oberfläche des Sonnenkörpers unter den Bedingungen des Gleichgewichts stärker von der Kugelgestalt abwiche, als es nach den Beobachtungen der Fall ist. Auch die beobachteten Merkurdurchgänge kommen hier in Betracht. Die daraus abgeleiteten Breiten des Merkur beruhen auf der Annahme, dass der Mittelpunkt der Sonnengestalt in der Ebene der Ekliptik selbst liegt. Daher würden alle beobachteten Breiten zu weit südlich liegen. Die daraus abgeleitete Länge des aufsteigenden Knotens würde zu groß, die des absteigenden zu klein ausfallen, und beide würden nicht übereinstimmen. Die Zahlen, die Sie in Ihrer Merkurtheorie mitgeteilt haben, deuten tatsächlich auf etwas Derartiges hin, scheinen sich aber mit einem so großen Unterschied, wie ihn die zwölfzolligen Kreise ergeben, nicht zu vertragen. Freilich gibt es zu wenige beobachtete Durchgänge durch den absteigenden Knoten. Da Sie keine Einzelheiten zu den ausgewählten Beobachtungen und zur Art der Ableitung der Breiten mitgeteilt haben, lässt sich nicht beurteilen, welche Genauigkeit man den einzelnen Breiten zuschreiben soll. Es wäre aber wohl der Mühe wert, die Merkurdurchgänge unter diesem Gesichtspunkt noch einmal aufs Sorgfältigste zu untersuchen.

\par\indent Es wäre außerordentlich interessant, wenn wir Beobachtungen der Sonne und der südlichen Zirkumpolarsterne mit Repetitionskreisen von der südlichen Erdhalbkugel hätten. Die südliche Breite müsste aber so groß sein, dass die Refraktion bei den Zirkumpolarsternen keine erhebliche Unsicherheit verursachen könnte.

\par\indent Bei dieser Gelegenheit bemerke ich noch, dass mir Herr Prof. \textsc{Schumacher} die Polhöhe mitgeteilt hat, die sein Bruder, Herr Kapitän \textsc{Schumacher}, in Friedrichsmark mit einem \textsc{Reichenbach}schen astronomischen Theodoliten beobachtet hat, der dem Ihren sehr ähnlich ist. Aus Beobachtungen des Polarsterns ergab sich $=55^\circ\,58'\,41''3$, aus Sonnenbeobachtungen $=55^\circ\,58'\,36''6$. Der Unterschied ist fast genau derselbe wie bei meinen Beobachtungen.

\Needspace{16\baselineskip}\par\indent\emph{Beobachtungen des Polarsterns in der unteren Kulmination an der neuen Göttinger Sternwarte.}

\begin{center}
\fontsize{9}{11}\selectfont
\setlength{\tabcolsep}{2pt}
\begin{tabular}{l|c|c|c|c|c|c}
1817 & \begin{tabular}{c}Anz. der\\Beob.\end{tabular}
& \begin{tabular}{c}Scheinbare\\Zenitdistanz\end{tabular}
& Barometer & \begin{tabular}{c}Thermo-\\meter\end{tabular}
& \begin{tabular}{c}Re-\\fraktion\end{tabular}
& \begin{tabular}{c}Refr. nach\\Bessel\end{tabular} \\
\hline
März 28 & 16 & $40^\circ\,7'\,22''40$ & $27^{\mathrm z}\,8^{\mathrm l}2$ & $-0^\circ5$ & $50''78$ & $50''25$ \\
April\ \ 1 & 22 & $24.98$ & $14.9$ & $+3.0$ & $50.88$ & \\
\ \ \ \ \ \ \ \ 2 & 24 & $23.38$ & $13.3$ & $+3.7$ & $50.49$ & \\
\ \ \ \ \ \ \ \ 8 & 28 & $26.36$ & $7\;2$ & $+4.4$ & $49.40$ & $48.87$ \\
\ \ \ \ \ \ \ 11 & 26 & $25.00$ & $9.3$ & $-1.4$ & $51.16$ & \\
\ \ \ \ \ \ \ 13 & 12 & $29.83$ & $8.5$ & $0$ & $50.68$ & \\
Mai\ \ \ 7 & 30 & $34.68$ & $10.9$ & $+7.5$ & $49.22$ & $48.67$ \\
\end{tabular}
\end{center}

\Needspace{12\baselineskip}\par\indent Daraus folgt

\begin{center}
\fontsize{9}{11}\selectfont
\setlength{\tabcolsep}{3pt}
\begin{tabular}{l|c|c|c}
1817 & Wahre Zenitdistanz & Scheinbare Polardistanz & Polhöhe \\
\hline
März 28 & $40^\circ\,8'\,13''18$ & $1^\circ\,40'\,2''51$ & $51^\circ\,31'\,49''33$ \\
April\ \ 1 & $15.86$ & $3.70$ & $47.84$ \\
\ \ \ \ \ \ \ \ 2 & $13.87$ & $4.00$ & $50.13$ \\
\ \ \ \ \ \ \ \ 8 & $15.76$ & $5.79$ & $50.03$ \\
\ \ \ \ \ \ \ 11 & $16.12$ & $6.97$ & $50.51$ \\
\ \ \ \ \ \ \ 13 & $20.51$ & $7.26$ & $46.75$ \\
Mai\ \ \ 7 & $23.90$ & $13.50$ & $49.67$ \\
\end{tabular}
\end{center}

\SourceLeaf{402}
\par\indent Also im Mittel aus 158 Beobachtungen $51^\circ\,31'\,49''42$, oder auf die Mitte der Sternwarte reduziert

\begin{center}\fontsize{10}{12}\selectfont $51^\circ\,31'\,49''30$ [$-0.644\,Q$ Handschriftliche Aufzeichnung.]\end{center}

\par\indent Die Refraktion ist den Tafeln von \textsc{Carlini} entnommen, die Polardistanz denen von \textsc{Bessel}. Meine Beobachtungen des Polarsterns in der unteren Kulmination von 1813 an der alten Sternwarte ergeben nach denselben Tafeln für deren Polhöhe $51^\circ\,31'\,55''48$. Auf die neue Sternwarte reduziert, ergeben die 136 Beobachtungen also

\[ 51^\circ\,31'\,49''37 \]

\par\indent Die Änderung des Gegengewichts hatte demnach überhaupt keinen Einfluss auf die Beobachtungen.

\par\indent Die obere Kulmination des Polarsterns konnte ich noch nicht beobachten. Am 19. Februar beobachtete ich 6 Zenitdistanzen in einem Abstand von 2 Stunden von der oberen Kulmination. Ihre Berechnung will ich hier noch als Beispiel beifügen, um mein Versprechen in der M.\,C. 1813, S. 484 zu erfüllen. Mit Vergnügen sehe ich aus dem letzten Heft Ihrer Zeitschrift, dass auch Herr \textsc{Littrow} diese Beobachtungsart mit Erfolg angewandt hat.

\Needspace{11\baselineskip}\par\indent Meine Beobachtungszeiten nach der Uhr waren

\begin{center}\fontsize{10}{12}\selectfont\setlength{\tabcolsep}{3pt}\begin{tabular}{ccc}
$2^{\mathrm u}$ & $54^{\mathrm m}$ & $19^{\mathrm s}5$\\
 & $59$ & $55.0$\\
$3$ & $4$ & $57.0$\\
 & $8$ & $36.0$\\
 & $13$ & $24.0$\\
 & $15$ & $31.0$\\
\end{tabular}\end{center}

\par\indent (Die großen Zeitabstände waren eine Folge des ungünstigen Wetters, das den ganzen Winter über nur für einzelne Augenblicke Sterne sichtbar werden ließ). Der durchlaufene Bogen betrug $222^\circ\,23'\,53''$, der Barometerstand $27^{\mathrm z}\,10^{\mathrm l}3$, der Thermometerstand $+4^\circ5$ und die Kulminationszeit des Sterns $0^{\mathrm u}\,55^{\mathrm m}\,23^{\mathrm s}59$. Mein Berechnungsverfahren will ich hier nur an seiner Anwendung erläutern.

\par\indent Das Mittel aller Uhrzeiten ist $3^{\mathrm u}\,6^{\mathrm m}\,7^{\mathrm s}08$. Die Unterschiede davon ohne Rücksicht auf das Vorzeichen sind

\begin{center}\fontsize{10}{12}\selectfont\setlength{\tabcolsep}{3pt}\begin{tabular}{cc}
$11^{\mathrm m}$ & $47^{\mathrm s}58$\\
$6$ & $12.08$\\
$1$ & $10.08$\\
$2$ & $48.92$\\
$7$ & $16.92$\\
$9$ & $23.92$\\
\end{tabular}\end{center}

\par\indent Mit diesen Unterschieden gehe ich in die bekannte Korrekturtafel der Höhen in Meridiannähe ein (ich benutze die von Herrn von \textsc{Zach} in der Attraction des montagnes) und erhalte

\begin{center}\fontsize{10}{12}\selectfont\setlength{\tabcolsep}{3pt}\begin{tabular}{c}
$273''01$\\
$75.50$\\
$2.68$\\
$15.56$\\
$104.11$\\
$173.42$\\
\end{tabular}\end{center}

\noindent Das Mittel daraus ist $107''38$. Ihm entspricht in derselben Tafel das Argument $7^{\mathrm m}\,23^{\mathrm s}73$. Daraus schließe ich (mithilfe eines kleinen Kunstgriffs, dessen Grund man sich leicht selbst klarmachen oder in den Götting. gel. Anz. 1815, S. 454 [23. März, siehe die weiter unten folgenden Anzeigen in diesem Band] \fancyhead[L]{\small SONNENBEOBACHTUNGEN.}\SourceLeaf{403} nachsehen kann), dass das Mittel der \emph{berechneten} Zenitdistanzen für die \emph{sechs} wirklichen Beobachtungszeiten annähernd gleich dem Mittel derjenigen für die \emph{zwei} Uhrzeiten ist

\begin{center}\fontsize{9}{11}\selectfont\setlength{\tabcolsep}{3pt}\begin{tabular}{ccccccccccc}
$3^{\mathrm u}$ & $6^{\mathrm m}$ & $7^{\mathrm s}08$ & $-$ & $7^{\mathrm m}$ & $23^{\mathrm s}73$ & $=$ & $2^{\mathrm u}$ & $58^{\mathrm m}$ & $43^{\mathrm s}35$\\
$3$ & $6$ & $7{,}08$ & $+$ & $7$ & $23{,}73$ & $=$ & $3$ & $13$ & $30{,}81$\\
\end{tabular}\end{center}

\Needspace{12\baselineskip}\noindent denen die Stundenwinkel $30^\circ\,49'\,56''9$ und $34^\circ\,31'\,48''3$ entsprechen. Die für diese Stundenwinkel nach bekannten Formeln berechneten Zenitdistanzen, mit der aus \textsc{Bessels} Tafel entnommenen Deklination $88^\circ\,20'\,7''53$ und der Polhöhe $51^\circ\,31'\,47''0$, sind

\begin{center}\fontsize{10}{12}\selectfont\begin{tabular}{lrrr}
& $37^\circ$ & $2'$ & $57''31$\\
& $37$ & $6$ & $32.75$\\\cline{2-4}
also das Mittel & $37$ & $4$ & $45.03$\\
\end{tabular}\end{center}
\noindent oder unter Berücksichtigung der mit $d\delta$, $d\varphi$ bezeichneten Korrekturen der zugrunde gelegten Deklination und Polhöhe\footnote{Die Koeffizienten sind mit dem mittleren Stundenwinkel einer von Herrn \textsc{Westphal} für die hiesige Polhöhe berechneten Tafel entnommen.}

\[ = 37^\circ\,4'\,45''03 - 1.00\,d\varphi + 0.83\,d\delta \]

\par\indent Das Mittel der beobachteten Zenitdistanzen ist $37^\circ\,3'\,58''83$ oder um die Refraktion $44''64$ bereinigt

\[ = 37^\circ\,4'\,43''47 \]

\par\indent Daraus folgt also $d\varphi = +1''56 + 0.83\,d\delta$ oder für die Polhöhe\newline \mbox{$= 51^\circ\,31'\,48''56 + 0.83\,d\delta$}. Die Beobachtungen in der unteren Kulmination haben $51^\circ\,31'\,49''42 - d\delta$ ergeben. Daraus würden $d\delta = +0''47$ und die Polhöhe $51^\circ\,31'\,48''95$ folgen. (Natürlich ist dies nur als Beispiel hinzugefügt. Um die Polhöhe von der zugrunde gelegten Deklination unabhängig zu machen, müssten die Beobachtungen in der Nähe der oberen Kulmination viel zahlreicher sein, und die Ergebnisse müssten dann nach der Methode der kleinsten Quadrate zusammengeführt werden). Man könnte auch eine Korrektur der zugrunde gelegten Rektaszension einführen. Man lässt dies aber besser weg, da dieses Element dank Ihrer Bemühungen und der unseres \textsc{Bessel} jetzt so genau bestimmt ist.

\Needspace{22\baselineskip}\par\indent Ich lasse nunmehr meine Sonnenbeobachtungen folgen.

\begin{center}\fontsize{9}{11}\selectfont\setlength{\tabcolsep}{3pt}\begin{tabular}{c|c|c|c|c|c}
$1817$ & \shortstack{durchlaufener\\Bogen} & \shortstack{Summe der\\Reduktion} & \shortstack{Anz. der\\Beob.} & Barometer & Thermometer\\
\hline$\text{Juni}\ 13$ & $283^\circ\,11'\,9''75$ & $11'\,28''90$ & $10$ & $27^{\mathrm z}\,6^{\mathrm l}6$ & $20^\circ8$\\
$16$ & $451\;21\;38.75$ & $52\,1.70$ & $16$ & $13.1$ & $13.0$\\
$17$ & $563\;46\;30.75$ & $77\,32.11$ & $20$ & $11.0$ & $14.3$\\
$18$ & $449\;59\;10.75$ & $25\,45.01$ & $16$ & $7.7$ & $16.4$\\
$19$ & $337\;8\;25.25$ & $12\,31.98$ & $12$ & $7.7$ & $19.3$\\
$20$ & $562\;18\;1.50$ & $61\;30\;15$ & $20$ & $7.4$ & $17.0$\\
$21$ & $337\;0\;25.50$ & $19\,20.42$ & $12$ & $7.5$ & $21.0$\\
$22$ & $280\;43\;13.25$ & $7\,55.97$ & $10$ & $8.3$ & $22.0$\\
$23$ & $337\;14\;3.75$ & $26\,26.97$ & $12$ & $8.1$ & $20.7$\\
$24$ & $449\;48\;49.25$ & $31\,1.55$ & $16$ & $8.1$ & $22.5$\\
$25$ & $281\;19\;53.25$ & $18\,40.57$ & $10$ & $8.1$ & $20.4$\\
$26$ & $394\;13\;39.25$ & $23\,47.90$ & $14$ & $7.5$ & $19.0$\\
$27$ & $620\;53\;26.75$ & $73\,42.54$ & $22$ & $5.8$ & $17.5$\\
$29$ & $452\;46\;31.50$ & $40\,2.88$ & $16$ & $8.9$ & $16.0$\\
$30$ & $226\;51\;23.75$ & $21\,7.83$ & $8$ & $7.6$ & $19.0$\\
\end{tabular}\end{center}

\SourceLeaf{404}
\Needspace{23\baselineskip}\par\indent Hieraus finde ich ferner

\begin{center}\fontsize{9}{11}\selectfont\setlength{\tabcolsep}{3pt}\begin{tabular}{c|c|c|c|c|c}
$1817$ & Refraktion & Parallaxe & Breite & \shortstack{Reduktion auf\\die Sonnenwende} & \shortstack{Zenitdistanz zur\\Sonnenwende}\\
\hline$\text{Juni}\ 13$ & $29''22$ & $4.01$ & $-0''11$ & $14'\,30''92$ & $28^\circ\,3'\,52''26$\\
$16$ & $30.73$ & $3.99$ & $+0.26$ & $5\,59.44$ & $48.63$\\
$17$ & $30.33$ & $3.99$ & $+0.34$ & $3\,58.31$ & $55.30$\\
$18$ & $29.67$ & $3.99$ & $+0.38$ & $2\,21.96$ & $53.46$\\
$19$ & $29.25$ & $3.99$ & $+0.39$ & $1\,10.43$ & $54.66$\\
$20$ & $29.17$ & $3.99$ & $+0.37$ & $0\,23.71$ & $51.41$\\
$21$ & $28.99$ & $3.99$ & $+0.30$ & $0\,1.81$ & $48.91$\\
$22$ & $28.91$ & $3.99$ & $+0.20$ & $0\,4.75$ & $52.10$\\
$23$ & $29.11$ & $3.98$ & $+0.08$ & $0\,32.50$ & $51.77$\\
$24$ & $28.89$ & $3.98$ & $-0.06$ & $1\,24.96$ & $51.62$\\
$25$ & $29.12$ & $3.99$ & $-0.22$ & $2\,42.22$ & $49.96$\\
$26$ & $29.35$ & $3.99$ & $-0.37$ & $4\,24.19$ & $51.61$\\
$27$ & $29.47$ & $3.99$ & $-0.52$ & $6\,30.75$ & $53.49$\\
$29$ & $30.05$ & $4.01$ & $-0.79$ & $11\,57.68$ & $51.86$\\
$30$ & $29.60$ & $4.01$ & $-0.88$ & $15\,17.82$ & $53.88$\\
\end{tabular}\end{center}

\par\indent Im Mittel aus 214 Beobachtungen

\[ = 28^\circ\,3'\,52''15 \]

\par\indent Meine Wahrscheinlichkeitstheorie ergibt für dieses Resultat den wahrscheinlichen Fehler $= \pm 0''32$. Nach den Mailänder Ephemeriden (denen auch Herr Kapitän \textsc{Schumacher} bei der Reduktion seiner oben erwähnten Beobachtungen die Deklinationen der Sonne entnahm) beträgt die scheinbare Schiefe der Ekliptik $= 23^\circ\,27'\,52''9$, also die Polhöhe $51^\circ\,31'\,45''05$ oder auf den Mittelpunkt reduziert

\begin{center}\fontsize{10}{12}\selectfont $50^\circ\,31'\,45''21$ [$+0.472\,Q$ Handschriftliche Aufzeichnung.]\end{center}

\noindent um $4''09$ kleiner als die aus dem Polarstern bestimmte Polhöhe. Mit \textsc{Bessels} Schiefe der Ekliptik wäre der Unterschied noch um anderthalb Sekunden größer.

\par\indent Ich muss noch bemerken, dass der Thermometerstand am Barometer nicht mit aufgezeichnet wurde. Bei so mäßigen Zenitdistanzen kann es nur einen sehr kleinen Bruchteil einer Sekunde ausmachen, wenn man die Reduktion auf eine Normaltemperatur oder auf die Temperatur der Luft im Freien vernachlässigt. Sollte jemand die Beobachtungen nach anderen Refraktionstafeln und unter Berücksichtigung einer solchen Refraktion berechnen wollen, so wird er wenig fehlgehen, wenn er die Temperatur des Quecksilbers bei den Sonnenbeobachtungen jeweils um $2^\circ$ niedriger, bei den Polarsternbeobachtungen um $2^\circ$ höher als die Temperatur im Freien ansetzt. Zu dieser Bemerkung veranlasst mich der an sich sehr berechtigte Hinweis von Herrn Prof. \textsc{Pasquich} im neuen Heft der Zeitschrift, das mir gerade ausgehändigt wurde. [Handschriftlich: $51^\circ\,30'\,48''26 - 0.966\,Q$ Wintersonnenwende 1817.]
